\IfFileExists{stacks-project.cls}{%
\documentclass{stacks-project}
}{%
\documentclass{amsart}
}

% For dealing with references we use the comment environment
\usepackage{verbatim}
\newenvironment{reference}{\comment}{\endcomment}
%\newenvironment{reference}{}{}
\newenvironment{slogan}{\comment}{\endcomment}
\newenvironment{history}{\comment}{\endcomment}

% For commutative diagrams we use Xy-pic
\usepackage[all]{xy}

% We use 2cell for 2-commutative diagrams.
\xyoption{2cell}
\UseAllTwocells

% We use multicol for the list of chapters between chapters
\usepackage{multicol}

% This is generally recommended for better output
\usepackage{lmodern}
\usepackage[T1]{fontenc}

% For cross-file-references
\usepackage{xr-hyper}

% Package for hypertext links:
\usepackage{hyperref}

% For any local file, say "hello.tex" you want to link to please
% use \externaldocument[hello-]{hello}
\externaldocument[introduction-]{introduction}
\externaldocument[conventions-]{conventions}
\externaldocument[sets-]{sets}
\externaldocument[categories-]{categories}
\externaldocument[topology-]{topology}
\externaldocument[sheaves-]{sheaves}
\externaldocument[sites-]{sites}
\externaldocument[stacks-]{stacks}
\externaldocument[fields-]{fields}
\externaldocument[algebra-]{algebra}
\externaldocument[brauer-]{brauer}
\externaldocument[homology-]{homology}
\externaldocument[derived-]{derived}
\externaldocument[simplicial-]{simplicial}
\externaldocument[more-algebra-]{more-algebra}
\externaldocument[smoothing-]{smoothing}
\externaldocument[modules-]{modules}
\externaldocument[sites-modules-]{sites-modules}
\externaldocument[injectives-]{injectives}
\externaldocument[cohomology-]{cohomology}
\externaldocument[sites-cohomology-]{sites-cohomology}
\externaldocument[dga-]{dga}
\externaldocument[dpa-]{dpa}
\externaldocument[sdga-]{sdga}
\externaldocument[hypercovering-]{hypercovering}
\externaldocument[schemes-]{schemes}
\externaldocument[constructions-]{constructions}
\externaldocument[properties-]{properties}
\externaldocument[morphisms-]{morphisms}
\externaldocument[coherent-]{coherent}
\externaldocument[divisors-]{divisors}
\externaldocument[limits-]{limits}
\externaldocument[varieties-]{varieties}
\externaldocument[topologies-]{topologies}
\externaldocument[descent-]{descent}
\externaldocument[perfect-]{perfect}
\externaldocument[more-morphisms-]{more-morphisms}
\externaldocument[flat-]{flat}
\externaldocument[groupoids-]{groupoids}
\externaldocument[more-groupoids-]{more-groupoids}
\externaldocument[etale-]{etale}
\externaldocument[chow-]{chow}
\externaldocument[intersection-]{intersection}
\externaldocument[pic-]{pic}
\externaldocument[weil-]{weil}
\externaldocument[adequate-]{adequate}
\externaldocument[dualizing-]{dualizing}
\externaldocument[duality-]{duality}
\externaldocument[discriminant-]{discriminant}
\externaldocument[derham-]{derham}
\externaldocument[local-cohomology-]{local-cohomology}
\externaldocument[algebraization-]{algebraization}
\externaldocument[curves-]{curves}
\externaldocument[resolve-]{resolve}
\externaldocument[models-]{models}
\externaldocument[functors-]{functors}
\externaldocument[equiv-]{equiv}
\externaldocument[pione-]{pione}
\externaldocument[etale-cohomology-]{etale-cohomology}
\externaldocument[proetale-]{proetale}
\externaldocument[relative-cycles-]{relative-cycles}
\externaldocument[more-etale-]{more-etale}
\externaldocument[trace-]{trace}
\externaldocument[crystalline-]{crystalline}
\externaldocument[spaces-]{spaces}
\externaldocument[spaces-properties-]{spaces-properties}
\externaldocument[spaces-morphisms-]{spaces-morphisms}
\externaldocument[decent-spaces-]{decent-spaces}
\externaldocument[spaces-cohomology-]{spaces-cohomology}
\externaldocument[spaces-limits-]{spaces-limits}
\externaldocument[spaces-divisors-]{spaces-divisors}
\externaldocument[spaces-over-fields-]{spaces-over-fields}
\externaldocument[spaces-topologies-]{spaces-topologies}
\externaldocument[spaces-descent-]{spaces-descent}
\externaldocument[spaces-perfect-]{spaces-perfect}
\externaldocument[spaces-more-morphisms-]{spaces-more-morphisms}
\externaldocument[spaces-flat-]{spaces-flat}
\externaldocument[spaces-groupoids-]{spaces-groupoids}
\externaldocument[spaces-more-groupoids-]{spaces-more-groupoids}
\externaldocument[bootstrap-]{bootstrap}
\externaldocument[spaces-pushouts-]{spaces-pushouts}
\externaldocument[spaces-chow-]{spaces-chow}
\externaldocument[groupoids-quotients-]{groupoids-quotients}
\externaldocument[spaces-more-cohomology-]{spaces-more-cohomology}
\externaldocument[spaces-simplicial-]{spaces-simplicial}
\externaldocument[spaces-duality-]{spaces-duality}
\externaldocument[formal-spaces-]{formal-spaces}
\externaldocument[restricted-]{restricted}
\externaldocument[spaces-resolve-]{spaces-resolve}
\externaldocument[formal-defos-]{formal-defos}
\externaldocument[defos-]{defos}
\externaldocument[cotangent-]{cotangent}
\externaldocument[examples-defos-]{examples-defos}
\externaldocument[algebraic-]{algebraic}
\externaldocument[examples-stacks-]{examples-stacks}
\externaldocument[stacks-sheaves-]{stacks-sheaves}
\externaldocument[criteria-]{criteria}
\externaldocument[artin-]{artin}
\externaldocument[quot-]{quot}
\externaldocument[stacks-properties-]{stacks-properties}
\externaldocument[stacks-morphisms-]{stacks-morphisms}
\externaldocument[stacks-limits-]{stacks-limits}
\externaldocument[stacks-cohomology-]{stacks-cohomology}
\externaldocument[stacks-perfect-]{stacks-perfect}
\externaldocument[stacks-introduction-]{stacks-introduction}
\externaldocument[stacks-more-morphisms-]{stacks-more-morphisms}
\externaldocument[stacks-geometry-]{stacks-geometry}
\externaldocument[moduli-]{moduli}
\externaldocument[moduli-curves-]{moduli-curves}
\externaldocument[examples-]{examples}
\externaldocument[exercises-]{exercises}
\externaldocument[guide-]{guide}
\externaldocument[desirables-]{desirables}
\externaldocument[coding-]{coding}
\externaldocument[obsolete-]{obsolete}
\externaldocument[fdl-]{fdl}
\externaldocument[index-]{index}

% Theorem environments.
%
\theoremstyle{plain}
\newtheorem{theorem}[subsection]{Theorem}
\newtheorem{proposition}[subsection]{Proposition}
\newtheorem{lemma}[subsection]{Lemma}

\theoremstyle{definition}
\newtheorem{definition}[subsection]{Definition}
\newtheorem{example}[subsection]{Example}
\newtheorem{exercise}[subsection]{Exercise}
\newtheorem{situation}[subsection]{Situation}

\theoremstyle{remark}
\newtheorem{remark}[subsection]{Remark}
\newtheorem{remarks}[subsection]{Remarks}

\numberwithin{equation}{subsection}

% Macros
%
\def\lim{\mathop{\mathrm{lim}}\nolimits}
\def\colim{\mathop{\mathrm{colim}}\nolimits}
\def\Spec{\mathop{\mathrm{Spec}}}
\def\Hom{\mathop{\mathrm{Hom}}\nolimits}
\def\Ext{\mathop{\mathrm{Ext}}\nolimits}
\def\SheafHom{\mathop{\mathcal{H}\!\mathit{om}}\nolimits}
\def\SheafExt{\mathop{\mathcal{E}\!\mathit{xt}}\nolimits}
\def\Sch{\mathit{Sch}}
\def\Mor{\mathop{\mathrm{Mor}}\nolimits}
\def\Ob{\mathop{\mathrm{Ob}}\nolimits}
\def\Sh{\mathop{\mathit{Sh}}\nolimits}
\def\NL{\mathop{N\!L}\nolimits}
\def\CH{\mathop{\mathrm{CH}}\nolimits}
\def\proetale{{pro\text{-}\acute{e}tale}}
\def\etale{{\acute{e}tale}}
\def\QCoh{\mathit{QCoh}}
\def\Ker{\mathop{\mathrm{Ker}}}
\def\Im{\mathop{\mathrm{Im}}}
\def\Coker{\mathop{\mathrm{Coker}}}
\def\Coim{\mathop{\mathrm{Coim}}}

% Boxtimes
%
\DeclareMathSymbol{\boxtimes}{\mathbin}{AMSa}{"02}

%
% Macros for moduli stacks/spaces
%
\def\QCohstack{\mathcal{QC}\!\mathit{oh}}
\def\Cohstack{\mathcal{C}\!\mathit{oh}}
\def\Spacesstack{\mathcal{S}\!\mathit{paces}}
\def\Quotfunctor{\mathrm{Quot}}
\def\Hilbfunctor{\mathrm{Hilb}}
\def\Curvesstack{\mathcal{C}\!\mathit{urves}}
\def\Polarizedstack{\mathcal{P}\!\mathit{olarized}}
\def\Complexesstack{\mathcal{C}\!\mathit{omplexes}}
% \Pic is the operator that assigns to X its picard group, usage \Pic(X)
% \Picardstack_{X/B} denotes the Picard stack of X over B
% \Picardfunctor_{X/B} denotes the Picard functor of X over B
\def\Pic{\mathop{\mathrm{Pic}}\nolimits}
\def\Picardstack{\mathcal{P}\!\mathit{ic}}
\def\Picardfunctor{\mathrm{Pic}}
\def\Deformationcategory{\mathcal{D}\!\mathit{ef}}
\usepackage{cleveref}
\crefname{originallemma}{original lemma}{original lemmas}
\Crefname{originallemma}{Original lemma}{Original lemmas}
\makeatletter
\newlabel{etale-cohomology-lemma-proper-base-change}{{0DDE}{source}{}{sourcepointer-1}{}}
\newlabel{etale-cohomology-lemma-proper-base-change@cref}{{[originallemma][0][]0DDE}{[1][0][]source}{}{}{}}
\makeatother
\crefname{originallemma}{original lemma}{original lemmas}
\Crefname{originallemma}{Original lemma}{Original lemmas}
\makeatletter
\newlabel{etale-cohomology-lemma-proper-base-change-mod-n}{{0F0C}{source}{}{sourcepointer-2}{}}
\newlabel{etale-cohomology-lemma-proper-base-change-mod-n@cref}{{[originallemma][0][]0F0C}{[1][0][]source}{}{}{}}
\makeatother
\crefname{originallemma}{original lemma}{original lemmas}
\Crefname{originallemma}{Original lemma}{Original lemmas}
\makeatletter
\newlabel{lemma-base-change-f-shriek-separated}{{0F55}{source}{}{sourcepointer-3}{}}
\newlabel{lemma-base-change-f-shriek-separated@cref}{{[originallemma][0][]0F55}{[1][0][]source}{}{}{}}
\makeatother
\crefname{originallemma}{original lemma}{original lemmas}
\Crefname{originallemma}{Original lemma}{Original lemmas}
\makeatletter
\newlabel{lemma-shriek-composition}{{0F7J}{source}{}{sourcepointer-4}{}}
\newlabel{lemma-shriek-composition@cref}{{[originallemma][0][]0F7J}{[1][0][]source}{}{}{}}
\makeatother
\crefname{originallemma}{original lemma}{original lemmas}
\Crefname{originallemma}{Original lemma}{Original lemmas}
\makeatletter
\newlabel{lemma-shriek-proper-and-open-base-change}{{0F7F}{source}{}{sourcepointer-5}{}}
\newlabel{lemma-shriek-proper-and-open-base-change@cref}{{[originallemma][0][]0F7F}{[1][0][]source}{}{}{}}
\makeatother
\crefname{originalreference}{original reference}{original references}
\Crefname{originalreference}{Original reference}{Original references}
\makeatletter
\newlabel{remark-f-shriek-base-change-composition}{{0F77}{source}{}{sourcepointer-6}{}}
\newlabel{remark-f-shriek-base-change-composition@cref}{{[originalreference][0][]0F77}{[1][0][]source}{}{}{}}
\makeatother
\crefname{originalreference}{original reference}{original references}
\Crefname{originalreference}{Original reference}{Original references}
\makeatletter
\newlabel{sites-cohomology-remark-compose-base-change}{{0E46}{source}{}{sourcepointer-7}{}}
\newlabel{sites-cohomology-remark-compose-base-change@cref}{{[originalreference][0][]0E46}{[1][0][]source}{}{}{}}
\makeatother
\title{Base change isomorphism for compactly supported derived pushforward along separated finite type morphisms}
\author{The Stacks Project Authors}
\date{stacks-project:more-etale.tex:local-snapshot-20261012; GNU FDL 1.2 or later}
\begin{document}
\maketitle
\setcounter{subsection}{1}
\noindent\textit{Editorial note.} All delivered mathematical text is the author's own, copied from the pinned Stacks Project chapter file more-etale.tex; no statement or proof step was rewritten, reordered or removed. The counted span is the author environment and its own proof environment. References inside the span are handed to the assembler as navigation stubs that carry the canonical tag of the cited tag only; whether each of them is load-bearing is recorded on the stub itself.\par\medskip
\noindent\textit{Editorial note.} Disclosure: the assembly argument is omitted verbatim -- 'A straightforward argument which we omit shows that both follow if we can show that the isomorphisms (1), (2), (3) are compatible with base change' (proof lines 3690-3691). What is omitted is the bookkeeping assembly from the three generating compatibilities to the two deliverables (independence of the compactification and stacking compatibility); the conclusion is stated precisely and every input is anchored: (1) by the Sites composing-base-change remark, (2) by the chapter remark on base change compatibility of the composition isomorphism, and (3) by lemma-shriek-proper-and-open-base-change. The compactification-independence part parallels the cofiltered argument spelled out in the proof of lemma-shriek-well-defined. No other step is omitted.\par\medskip
\noindent\textit{Allowed standard prerequisites (editorial).}\begin{itemize}
\item Proper base change (Etale Cohomology, Lemmas etale-cohomology-lemma-proper-base-change 0DDE and etale-cohomology-lemma-proper-base-change-mod-n 0F0C). cited at proof lines 3675-3680. The first arrow of the construction: base change for the compactification pushforward, in the bounded below torsion and torsion ring cases respectively.
\item Base change for separated lower shriek (Lemma lemma-base-change-f-shriek-separated, canonical tag 0F55). cited at proof lines 3680-3681. The second arrow of the construction: the isomorphism for the lower shriek part, valid before sheafification issues since it is the abelian-sheaf level iso.
\item Base change compatibility of the exchange map (Lemma lemma-shriek-proper-and-open-base-change, canonical tag 0F7F). cited at proof line 3704. Anchors compatibility (3), the exchange-map input of the assembly argument.
\end{itemize}\medskip
\ifdefined\hypertarget\hypertarget{sourcepointer-1}{}\fi
\noindent\textit{Source reference.} Original lemma 0DDE: sources/stacks-chow/tags.txt, source lines 16764--16764. Cited at proof lines 3675-3677: the first arrow of the construction in the bounded below torsion case.\par\medskip
\ifdefined\hypertarget\hypertarget{sourcepointer-2}{}\fi
\noindent\textit{Source reference.} Original lemma 0F0C: sources/stacks-chow/tags.txt, source lines 18757--18757. Cited at proof lines 3678-3680: the first arrow of the construction when Lambda is a torsion ring.\par\medskip
\ifdefined\hypertarget\hypertarget{sourcepointer-3}{}\fi
\noindent\textit{Source reference.} Original lemma 0F55: sources/stacks-chow/tags.txt, source lines 18925--18925. Cited at proof lines 3680-3681: the second arrow of the construction.\par\medskip
\ifdefined\hypertarget\hypertarget{sourcepointer-4}{}\fi
\noindent\textit{Source reference.} Original lemma 0F7J: sources/stacks-chow/tags.txt, source lines 19009--19009. Cited at statement lines 3659-3660 and proof line 3689: the composition isomorphisms with which the base change isomorphisms must be compatible.\par\medskip
\ifdefined\hypertarget\hypertarget{sourcepointer-5}{}\fi
\noindent\textit{Source reference.} Original lemma 0F7F: sources/stacks-chow/tags.txt, source lines 19005--19005. Cited at proof line 3704: anchors compatibility (3) of the exchange map with base change.\par\medskip
\ifdefined\hypertarget\hypertarget{sourcepointer-6}{}\fi
\noindent\textit{Source reference.} Original reference 0F77: sources/stacks-chow/tags.txt, source lines 18997--18997. Cited at proof line 3703: anchors compatibility (2) of the lower shriek composition isomorphism with base change.\par\medskip
\ifdefined\hypertarget\hypertarget{sourcepointer-7}{}\fi
\noindent\textit{Source reference.} Original reference 0E46: sources/stacks-chow/tags.txt, source lines 17666--17666. Cited at proof lines 3701-3702: anchors compatibility (1) of the composition of proper pushforwards with base change.\par\medskip
\begin{lemma}
\label{lemma-base-change-shriek}
Consider a cartesian square
$$
\xymatrix{
X' \ar[r]_{g'} \ar[d]_{f'} & X \ar[d]^f \\
Y' \ar[r]^g & Y
}
$$
of quasi-compact and quasi-separated schemes with $f$ separated and
of finite type. Then there is a canonical isomorphism
$$
g^{-1} \circ Rf_! \to Rf'_! \circ (g')^{-1}
$$
Moreover, these isomorphisms are compatible with the isomorphisms
of Lemma \ref{lemma-shriek-composition}.
\end{lemma}

\begin{proof}
Choose a compactification $j : X \to \overline{X}$ over $Y$
and denote $\overline{f} : \overline{X} \to Y$ the structure morphism.
Let $j' : X' \to \overline{X}'$ and $\overline{f}' : \overline{X}' \to Y'$
denote the base changes of $j$ and $\overline{f}$.
Since $Rf_! = R\overline{f}_* \circ j_!$ and $Rf'_! =
R\overline{f}'_* \circ j'_!$ the isomorphism can be constructed via
$$
g^{-1} \circ R\overline{f}_* \circ j_! \to
R\overline{f}'_* \circ (\overline{g}')^{-1} \circ j_! \to
R\overline{f}'_* \circ j'_! \circ (g')^{-1} 
$$
where the first arrow is the isomorphism given to us by the
proper base change theorem (\'Etale Cohomology, Lemma
\ref{etale-cohomology-lemma-proper-base-change} in the bounded below torsion
case and \'Etale Cohomology, Lemma
\ref{etale-cohomology-lemma-proper-base-change-mod-n} in the
case that $\Lambda$ is torsion) and the second arrow is the isomorphism of
Lemma \ref{lemma-base-change-f-shriek-separated}.

\medskip\noindent
To finish the proof we have to show two things: first we have to show
that the isomorphism of functors so obtained does not depend on
the choice of the compactification and second we have to show that
if we vertically stack two base change diagrams as in the lemma, then
these base change isomorphisms are compatible with the isomorphisms
of Lemma \ref{lemma-shriek-composition}.
A straightforward argument which we omit shows that both follow
if we can show that the isomorphisms
\begin{enumerate}
\item $Rg_* \circ Rf_* = R(g \circ f)_*$ for $f : X \to Y$ and $g : Y \to Z$
proper,
\item $g_! \circ f_! = (g \circ f)_!$ for $f : X \to Y$ and $g : Y \to Z$
separated and quasi-finite, and
\item $g_! \circ Rf'_* = Rf_* \circ g'_!$ for  $f : X \to Y$ and
$f' : X' \to Y'$ proper and $g : Y' \to Y$ and $g' : X' \to X$
separated and quasi-finite with $f \circ g' = g \circ f'$
\end{enumerate}
are compatible with base change. This holds for (1) by
Cohomology on Sites, Remark \ref{sites-cohomology-remark-compose-base-change},
for (2) by Remark \ref{remark-f-shriek-base-change-composition}, and
(3) by Lemma \ref{lemma-shriek-proper-and-open-base-change}.
\end{proof}

\end{document}
